\section{Introduction}
%\note{
%Answer the following questions
%\begin{itemize}
%  \item What is the problem you are solving?
%  \begin{itemize}
%    \item Why is it important?
%    \item How prevalent is it?
%  \end{itemize}
%  \item How is it solved currently, and what are the limitations of current methods?
%  \item What is your approach? How does it overcome the limitations above?
%  \item How do you know your approach works? (Evaluation)
%  \item What difference does it make?
%\end{itemize}
%
%}

Digital traces of human activity have exposed social behavior to web and data mining algorithms. Researchers have analyzed the growing social data to understand, among other things, how information spreads in social networks~\cite{Romero11www}, and the factors affecting user engagement~\cite{Barbosa2016} and performance~\cite{Singer2016plosone} in online platforms. Yet social data analysis presents multi-faceted challenges that existing data mining approaches are not well-equipped to handle.
Real-world data is noisy and sparse. To uncover hidden patterns, scientists aggregate data over the population, which tends to improve statistics and signal-to-noise ratio. However, real-world data is also heterogeneous, i.e., composed of subgroups that vary widely in size and behavior. This heterogeneity can severely bias analysis of social data. %estimates of causal effects.

One example of such bias is Simpson's paradox. According to the paradox, %which is often encountered in economics, demographics and clinical research,
an association observed in data that has been aggregated over an entire population may be quite different from, and even opposite to, those of the underlying subgroups. Failure to take this effect into account can distort conclusions drawn from data.
One of the better known examples of Simpson's paradox comes from a study of gender bias in graduate admissions~\cite{Bickel1975}. In the aggregate admissions data there appears to be a statistically significant bias against women: a smaller fraction of female applicants is admitted for graduate studies. However, when admissions data is disaggregated by department, women have parity and even a slight advantage over men in some departments. The paradox arises because departments preferred by female applicants have lower admissions rates for both genders.


Simpson's paradox also affects analysis of trends. When measuring how an outcome changes as a function of an independent variable, the characteristics of the population over which the trend is measured may change as a function of the independent variable due to survivor bias. As a result, the data may appear to exhibit a trend, which disappears or reverses when the data is disaggregated. %The trends discovered in disaggregated are more robust and more likely to describe individual behavior than the trends found in aggregate data.
For example, it may appear that repeated exposures to information or hashtags on social media make an individual less likely to share it with his or her followers~\cite{Romero11www}. In fact, the opposite is true: the more people are exposed to information, the more likely they are to share it with followers~\cite{Lerman2016futureinternet}. However, those who follow many others (and are likely to be exposed to a meme or a hashtag multiple times) are less responsive overall, due to the high volume of information they receive~\cite{Hodas12socialcom}. Their reduced susceptibility to exposures biases the aggregate response, leading to wrong conclusions about behavior. Once disaggregated based on the volume of information received, a clearer pattern of exposure response emerges, one that is more predictive of the actual response~\cite{Hodas14srep}.
However, despite accumulating evidence that Simpson's paradox %could distort the associations found in
affects inference of trends in social and behavioral data~\cite{Singer2016plosone,Barbosa2016,KootiA2017www,Agarwal2017quitting}, researchers do not routinely test for it in their studies.

% What we are doing
We describe a method to identify Simpson's paradoxes in the analysis of trends in social data. %Our method identifies trends in the aggregate data that disappear or reverse when the data is disaggregated into its underlying subgroups.
%Our method automatically identifies covariates in data that could be used to disaggregate it into more homogeneous subgroups.
Our statistical approach finds pairs of variables---that we call \emph{Simpson's pairs}---such that a trend in some outcome as a function of the first \emph{independent} variable observed in the aggregate data  disappears or reverses itself when the data is disaggregated into distinct subgroups on the second \emph{conditioning} variable.
We perform mathematical analysis, which identifies two necessary conditions for the paradox to occur: (1) the independent and conditioning variables are correlated and (2) the value of the outcome variable differs within conditioning subgroups.
%\begin{itemize}
%  \item The distribution of the conditioning variable on which the data is disaggregated changes as a function of the independent variable; i.e., the composition of the subgroups depends on the independent variable
%
%\item For a given value of the independent variable, the expected value of the outcome
%%cannot be constant with respect to the disaggregation variable.
%is different in different subgroups.
%\end{itemize}

% Evaluation
We apply the proposed approach to data collected from Stack Exchange, a popular question-answering platform. Specifically, we analyze factors affecting answerer performance, measured as the likelihood the answer provided by the answerer will be accepted by the asker as the best answer to his or her question. We construct a variety of features to describe answers and answerers, and study how the outcome---in this case answer acceptance---depends on these features. We compare results of statistical trends found in aggregated and disaggregated data to identify Simpsons's pairs. Our analysis discovers several cases of Simpson's paradox in Stack Exchange data. In addition to a known effect that describes deterioration of answerer performance over the course of a session, our method identifies several new cases of Simpson's paradox. These paradoxes yield new insights into answerer performance on Stack Exchange. For example, creating a new feature to describe answerers, which combines variables of a Simpson's pair, leads to a more robust proxy of answerer performance.

% Summary
%Social data often comes from populations comprised of subgroups that differ in their behavior. 
Presence of a Simpson's paradox in social data can indicate interesting patterns~\cite{fabris2000discovering}, including important behavioral differences. The paradox suggests that subgroups within the population under study systematically differ in their behavior, and these differences are large enough to affect aggregate trends. In such cases the trends discovered in disaggregated data are more likely to describe---and predict---individual behavior than the trends found in aggregated data. Thus, to build more robust models of behavior, data scientists need to identify the subgroups within their data %whose behavior could affect their conclusions.
or risk drawing wrong conclusions.
The method presented in this paper provides a simple framework for identifying such interesting subgroups by systematically searching for Simpson's paradox in trend data.
%Simpson's paradox as a tool to discover new important features of data. e.g., we discover a better measure of answerer quality.


The rest of the paper is organized as follows. First, we present background, including multiple observations of Simpson's paradox in a variety of disciplies (Sec.~\ref{sec:related}). Then we describe our methodology for detecting Simpson's paradox by identifying covariates in data, and analyze the paradox mathematically to gain more insight into its origins (Sec.~\ref{sec:methods}). Finally, we apply our method to real-world data from the question-answering site Stack Exchange, and demonstrate its ability to automatically identify novel cases of Simpson's paradox (Sec.~\ref{sec:results}). We conclude with the discussion of implications.

